\paragraph{Integral sections for the height comparison.}
\label{par:selected-height-sections}
\hypertarget{source-selected-height-sections}{}
\begin{established}[Complementary weights and the gaps of the selected rows]
\label{hyp:height-section-weights}\label{cor:selected-height-sections-1}
Fix the coefficient field, the two level primes and either orientation.
Write \(u_0\) for the weight with zero row entries and central
coordinate one in the group of weights. The auxiliary section before
the last lift has weight \(i_\eta g_0^*u_0\). Put
\[
 w_1=2(w_L-u_0),\qquad w_2=(s+4)w_L-w_F,\qquad
 d=s+2,\qquad T'=T+d .
\]
The two central degrees are \(2n_W-2\) and \((s+4)n_W-n_P\).
The run satisfies
\[
 8fr+4=n_P\le n_W,\qquad (s+4)\operatorname{DF}(w_L)+2
       \le\min(g_0,g_1),\qquad s\ge3,\qquad T\ge40d .
\]
Every interior adjacent difference \(a\) in a nonempty block of
\(\kappa\) satisfies \(a+2\ge T'(T'+2r)\). Each terminal entry \(a\)
satisfies \(a\ge2T^2\); its offset is \(d_0=dn_W+2r\) in the
positive block and zero in the negative block, with \(0\le d_0\le T\).
\end{established}
\begin{proof}
The central degrees follow by addition of weights.
Lemma~\ref{hyp:common-slice-scales} gives the degree and order bounds.
Since \(D_L=\operatorname{DF}(w_L)\le 2frb\), \(s,f,r,k\) are
polynomially bounded in \(x\), and the two primes satisfy
\(g_i\ge(kb)^8\), the level inequality follows at all large scales.

We give the row argument, including the margin needed at a block end.
Let \(h=\eta_{\rm size}/m\) be the spacing of the central parameters,
and \(\Delta\) the separation parameter of the outer ones. In the
successive choices of Paragraph~\ref{par:window-successive-bounds}
one can impose
\[
       h\ge4T^2+2t,\qquad
       \Delta\ge(T+D_*+2)(T+D_*+2+2r).
\]
Indeed these lower bounds are fixed polynomials in \(x\) after
Lemma~\ref{lem:rank-window}. In choosing the width and separation
before rounding the nodes, adjoining either polynomial to their lower
bounds costs \(O(\log x)\) on the logarithmic scale; the available
upper bounds there are \(x^{\epsilon/4}\) for \(\log\eta_{\rm size}\)
and \(x^{\epsilon/4}/2\) for \(\log\Delta\). Equivalently the
degree-power lower bounds suffice with exponent three. For
\(N=k+m(t+1)\), Lemma~\ref{hyp:common-slice-scales} gives
\[
    N\ge t\ge T,\qquad
    (T+D_*+2)(T+D_*+2+2r)\le(41T/40)(81T/40)\le N^3 .
\]
These choices give \(G:=T'(T'+2r)\le\min(\Delta,h)\), since
\(d\le T/40\) and \(2r\le T\).

Fix an embedding and number its central nodes \(b_0,\ldots,b_{m-1}\).
Their spacings and the corresponding integral centres \(e_i\) satisfy
\[
     h\le b_{i+1}-b_i\le h+1,\qquad |e_i-b_i-h/2|\le2 .
\]

The two errors come from rounding the nodes and their half-shifts.
The decreasing list for the positive block consists of outer
parameters \(v_0,\ldots,v_{k-2}\) followed by
\(-b_0,\ldots,-b_{j-1}\). Within the outer list consecutive differences
are at least \(\Delta\), by its spacing, and the central differences
are at least \(h\). At the junction \(v_{k-2}\ge0\); the lower endpoint
of the translated windows gives \(e_0\ge t\eta_{\rm size}\). Thus
\[
   v_{k-2}-(-b_0)\ge b_0\ge e_0-\eta_{\rm size}
                   \ge(t-1)\eta_{\rm size}\ge h .
\] In the reversed negative block the consecutive
unshifted differences are again differences of adjacent \(b_i\)'s.
Subtracting the row shift changes an interior difference by one; in
particular \(a+2\ge G\) in either block.

At lift \(1\le j\le m\), substitution of the last row in the two
weight formulas gives respectively
\[
 e_{j-1}-b_{j-1}-(t+1)/2,\qquad
 b_j-e_{j-1}-(t+1)/2 \quad(j<m).
\]
The centre error and the lower bound on node spacing make both at
least \((h-4-t-1)/2\ge2T^2\), since \(t\ge5\).
The positive offset is at most \(T\) by
Lemma~\ref{hyp:common-slice-scales}; the empty last block at \(j=m\)
requires no terminal inequality.
\end{proof}

\begin{proposition}[Bounded complementary sections at the selected weights]\label{prop:height-complementary-sections}
For the run of Lemma~\ref{hyp:height-section-weights}, the weights \(w_1,w_2\) used at the
corresponding lifts are admissible. There are integral
sections \(t_1,t_2\) at these weights, nonzero on every
component, with
\[
  \tfrac12\log\|t_1\|_{\sup,\sigma}\le P_4 sS_L,\qquad
              \log\|t_2\|_{\sup,\sigma}\le P_4 sS_L
\]
with \(P_4=\max\{1,5(1009Z_b+\lambda_*)+162 Z_b^2/s\}\),
where \(Z_b=frL\) is the parameter of Proposition~\ref{prop:slice-evaluation-net} and
\(\lambda_*\) bounds the logarithmic cost of the integralizing scalar
per unit of weight size. The choices use the fixed
coefficient field and light columns. At the definite lift
only \(t_2\) is needed.
\end{proposition}

\begin{proof}
\proofstep{The allowable row-entry shifts.}
Here is the row-entry check. Put \(d=s+2\) and \(T'=T+d\).
At an interior gap in either sign block let \(a\) be the gap of
\(\kappa\), and \(l=\lfloor(a+2)/T'\rfloor\) the corresponding
gap of \(w_L\). The page gives \(a+2\ge T'(T'+2r)\). Write
\(a+2=T'l+\rho\) with \(0\le\rho<T'\). The gap of \(w_2\) is
\[
 (T'+2)l-a=2l-\rho+2>0.
\] At a last entry \(a\ge2T^2\) of
\(\kappa\), put \(d_0=dn_W+2r\) in the positive block and
\(d_0=0\) in the negative block.

Now
\(l=\lfloor(a-d_0)/T'\rfloor\) and \(0\le d_0\le T\).
Because \(T\ge40d\), we have
\[
 T'\left(1+\tfrac32T\right)\le 2T^2-T .
\]
Write \(a-d_0=T'l+\rho\) with \(0\le\rho<T'\). The last entry
of \(w_2\) in this block is positive, since
\[
 l=\left\lfloor\frac{a-d_0}{T'}\right\rfloor\ge\frac{3T}{2},
 \qquad 2l-d_0-\rho>3T-T-T'=T-d>0 .
\]
The rows of \(w_1\) are twice the rows of \(w_L\).
The row and central bounds for \(w_L\), with \(2,s+4\le T\),
give \(w_1,w_2\le\kappa\) componentwise.
For either weight \(u\), the rank and level tests follow from
\[
 \begin{aligned}
 n(u)&\in\{2n_W-2,(s+4)n_W-n_P\}, &n(u)&\ge n_W\ge r,\\
 \operatorname{DF}(u)+2&\le (s+4)\operatorname{DF}(w_L)+2,
 &\operatorname{DF}(u)+2&\le\min(g_0,g_1).
 \end{aligned}
\]
In particular 
\[
  g_0g_1\ge(\operatorname{DF}(u)+2)^2
\ge4\operatorname{DF}(u)+4.
\]
 The size bounds are
\(\lvert w_1\rvert\le2S_L\) and \(\lvert w_2\rvert\le(s+4)S_L\).
\label{cor:selected-height-sections-2}

\proofstep{The two sections from the nonvanishing jet test.}
Fix \(u\in\{w_1,w_2\}\) among the weights needed at this lift.
Work at a good complex embedding. On each component \(a\) choose an
orthonormal frame. The positive last entries and \(n(u)\ge r\)
allow the block determinant harmonic polynomial. Let \(D_u\) denote its nonzero leading differential operator.
At the level just checked, Lemma~\ref{lem:ideal-lattice-jet} gives
\[
 d=\operatorname{ord}(D_u)\le\operatorname{DF}(u),\qquad
        D_u\theta_\nu(0)\ne0\quad\text{for some }\nu .
\]
The corresponding theta evaluation is linear in the polynomial.
The bounded grid of determinant polynomials in the proof of
Proposition~\ref{prop:three-shifted-spanning-sets} spans the same
harmonic space. Fix a translate \(\nu\) on which the jet is nonzero. If every grid
polynomial vanished on this translate, linearity would give zero for the
original polynomial there as well. Choose one grid rule \(F_a\) with
nonzero value at this frame and translate.

Let \(A_u\) be the sum of the flag entries and put \(h_u=frn(u)\).
The grid estimate and the complex theta norm calculation in that
proposition use only \(d\le A_u\), \(r\le n(u)\) and the logarithmic
size bounds at \(u\). These hold by the coordinatewise bounds
\(u\le\kappa\) above. Before applying the integralizing scalar they give
at every embedding \(\sigma\)
\[
 \log\|F_a\|_{\sup,\sigma}
 \le 140rL A_u+10L A_u+200h_uL
 \le 400 Z_b |u| .
\]
Let \(x_a\) be the chosen frame on component \(a\), and put
\(C=\{0,\ldots,M_c\}^{M_c}\).

For fixed \(m_b\) with \(b\ne a\),
the coefficient of \(m_a\) in \(\sum_b m_bF_b(x_a)\) is nonzero.
Consequently
\begin{equation}\label{eq:height-section-grid-avoidance}
 \begin{aligned}
 H_a&=\{m\in C:\textstyle\sum_b m_bF_b(x_a)=0\},&
       |H_a|&\le(M_c+1)^{M_c-1},\\
 \left|C\setminus\bigcup_a H_a\right|
     &\ge (M_c+1)^{M_c-1}.&&
 \end{aligned}
\end{equation}
By \eqref{eq:height-section-grid-avoidance} we can choose one \(m\)
avoiding every zero locus and multiply \(\sum_b m_b F_b\)
by the integralizing scalar \(s(u)\) of the theta module.
This scalar is nonzero and has log
absolute value at most \(\lambda_*|u|\) at every embedding.
The scaled combination lies in the integral cusp lattice at \(u\),
is nonzero on each analytic component and therefore on each
algebraic component. Its log supremum is at most
\[
               (400Z_b+\lambda_*)|u|+2\log M_c .
\]

For this estimate \(M_c\ge1\) and \(\sum_a m_a\le M_c^2\).
The marked-lattice count in Proposition~\ref{prop:slice-evaluation-net}
gives \(\log M_c\le41Z_b^2\), hence also \(\log M_c\le81Z_b^2\).
For \(|w_1|\le2S_L\), \(|w_2|\le(s+4)S_L\) and \(s\ge3\), put
\[
       P_4=\max\{1,5(1009Z_b+\lambda_*)+162 Z_b^2/s\}
\]
Since \(s+4\le5s\) and \(S_L\ge1\), the preceding bounds give
\[
 \tfrac12\log\|t_1\|_{\sup,\sigma}\le P_4 s S_L,\qquad
           \log\|t_2\|_{\sup,\sigma}\le P_4 s S_L .
\]
The two selections take place over the original coefficient field.
Their restrictions to an incidence field are nonzero because
that field contains the corresponding component field. The
selection uses only the section data of this run. In particular
\(t_2\) serves both the thickened weight-\(w_F\) line and the
comparison of its height with the light height.
\label{cor:selected-height-sections-3}
\end{proof}
